\documentclass[10pt,a4paper]{amsart}
\usepackage[margin=19mm]{geometry}
\usepackage{amsmath,amssymb,mathtools}
\usepackage[scr=rsfs,cal=euler]{mathalfa}
\usepackage{xfrac}
\usepackage[unicode,hidelinks,bookmarks=false]{hyperref}
\newcommand{\ie}{i.e.\ }
\newcommand{\cf}{cf.\ }
\newcommand{\resp}{respectively\ }
\newcommand{\Subsection}{\subsection}
\providecommand{\nobreakdash}{\nobreak\hbox{-}}
\setlength{\emergencystretch}{2em}
\multlinegap0pt
\usepackage[shortlabels]{enumitem}
\usepackage{tcolorbox}
\usepackage{bm}
\usepackage{xfrac}
\newtheorem{assumptions}{Assumptions}
\newtheorem{lemma}{Lemma}
\newcommand{\Lip}{\mathrm{Lip}}
\def\Lp#1{\mathrm{L}^{#1}}
\newcommand{\bbN}{\mathbb{N}}
\newcommand{\bbR}{\mathbb{R}}
\def\dd{\mathrm{d}}
\title{Multiple Neural Operators Achieve Near-Optimal Rates for Multi-Task Learning}
\author{Adrien Weihs, Hayden Schaeffer}
\date{Source: arXiv:2605.22724v1 (CC BY 4.0)}
\begin{document}
\maketitle

\noindent\emph{Source and licence.} Multiple Neural Operators Achieve Near-Optimal Rates for Multi-Task Learning. Version arXiv:2605.22724v1 (CC BY 4.0), distributed by arXiv under the Creative Commons Attribution 4.0 International licence, http://creativecommons.org/licenses/by/4.0/, as recorded in the arXiv licence metadata for this exact version and shown on the abstract page https://arxiv.org/abs/2605.22724v1. Full licence notice: \texttt{licenses/arxiv-2605.22724-CC-BY-4.0.txt}.\par\smallskip

\noindent\textit{Editorial scope.} Counted unit: the article's lemma lem:boundedLipCube (source0.tex lines 1158--1166). The article's complete original proof of it is delivered with it (source0.tex lines 1945--2034, one \texttt{proof} environment of 90 lines, which the article places away from the statement). No mathematics is written, repaired or re-derived: the statement and the proof are the article's own lines, in the article's own notation, and no retained line is substituted or re-flowed.  Retained uncounted as the source-local context the proof needs: lines 385--395; lines 1903--1911.  Nothing was omitted from any retained span, so the package carries no omission record.  No retained line contains a citation, so the wrapper carries no bibliography and no bibliography entry.  No retained line contains a figure environment, an image-inclusion command or drawing code, so no image is delivered and the package declares no asset dependency.  Editorial formatting only, in addition: an AMS \texttt{amsart} preamble that restates the article's own declarations and macros used by the retained lines, namely the environments \texttt{assumptions}, \texttt{lemma} on separate counters, together with the standard packages those lines need.  The retained environments print in this document as Assumptions 1, Lemma 1 in order of appearance, whereas the article prints the same items with the numbering of its own full document.  The complete native source of the article as distributed in the arXiv e-print for this exact version is bundled unchanged as \texttt{sources/201155/source0.tex}.
\begin{assumptions} We make the following assumption on our spaces.
\begin{enumerate}[label=\textbf{S.\arabic*}]
\item The space $U(d_U,\gamma_U,L_U,\beta_U)$ is a function set such that \begin{enumerate}
    \item any function $u \in U$ is defined on $\Omega_U := [-\gamma_U,\gamma_U]^{d_U}$;
    \item for all functions $u \in U$ and $x,y \in \Omega_U$, we have \[
    \vert u(x) - u(y) \vert \leq L_U \vert x - y \vert;
    \]
    \item for all functions $u \in U$, we have $\Vert u \Vert_{\Lp{\infty}} \leq \beta_U$.
\end{enumerate}\label{assumption:Main:assumptions:S4}
\end{enumerate}
\end{assumptions}

\begin{lemma}[$\Lp{r}$-norm of $\sin$] \label{lem:sin}
For every $1 \leq r < \infty$ and every nonzero multi-index $\kappa \in \mathbb N^d$, one has
\[
\|\sin(\kappa \cdot x)\|_{\Lp{r}([0,2\pi]^d)}
=
\left((2\pi)^{d-1}\int_0^{2\pi} |\sin t|^r \,\dd t\right)^{1/r}.
\]
In particular, the quantity $c_r := \|\sin(\kappa \cdot x)\|_{\Lp{r}([0,2\pi]^d)}$ is independent of $\kappa$.
\end{lemma}

\begin{lemma}[Infinite-dimensional Cube in Bounded Lipschitz Classes]
\label{lem:boundedLipCube}
Let $d_U \in \bbN$, $\gamma_U,L_U,\beta_U > 0$ and $U = U(d_U,\gamma_U,L_U,\beta_U)$ satisfy Assumption \ref{assumption:Main:assumptions:S4}. 
Then, viewed as a subset of the Banach space $\Lp{r}(\Omega_U)$ with $r\geq 1$, 
$U$ contains an infinite-dimensional hypercube $Q_\eta$ for every
\(
\eta > 1+\frac{1}{d_U}.
\)
\end{lemma}

\begin{proof}[Proof of Lemma \ref{lem:boundedLipCube}]
In the proof $C>0$ denotes a constant that may change from line to line, independently of the summation indices and $\kappa$.

Up to an affine rescaling of the domain, it suffices to treat the case
\(
\Omega_U=[0,2\pi]^{d_U}.
\)
For each multi-index $\kappa\in \mathbb N^{d_U}$, define
\[
\tilde e_\kappa(x) := \sin(\kappa\cdot x),
\qquad
e_\kappa(x):=\frac{\tilde e_\kappa(x)}{\|\tilde e_\kappa\|_{\Lp{r}(\Omega_U)}} = \frac{\tilde e_\kappa(x)}{c_r}
\]
where we used Lemma \ref{lem:sin} for the last equality. 
These functions are orthogonal in $\Lp{2}(\Omega_U)$, and hence linearly independent in $\Lp{r}(\Omega_U)$, with
\(
\|e_\kappa\|_{\Lp{r}(\Omega_U)} =1.
\)
For each $\kappa \in \bbN^{d_U}$, define $e^*_\kappa \in \Lp{r}(\Omega_U)^*$ as
\[
e_\kappa^*(u)
:=
\frac{2 c_r}{(2\pi)^d}\int_{[0,2\pi]^d} u(x)\sin(\kappa\cdot x)\, \dd x.
\]
By the identity $\sin(a)\sin(b) = \frac{1}{2}(\cos(a-b) - \cos(a+b))$, it is straightforward to check that $e_\kappa^*(e_{\kappa'})=\delta_{\kappa\kappa'}$. Moreover, by H\"older's inequality with $r'$ such that $1/r + 1/r' = 1$, 
\[
|e_\kappa^*(u)|
\le
\frac{2 c_r c_{r'}}{(2\pi)^d}
\|u\|_{\Lp{r}(\Omega_U)}
\le C \|u\|_{\Lp{r}(\Omega_U)},
\]
so
\(
\|e_\kappa^*\|_{\Lp{r}(\Omega_U)^*} \le C
\)
uniformly in $\kappa$.

Next, we enumerate $\{e_\kappa\}_{\kappa \in \bbN^{d_U}}$ as $\{e_j\}_{j=1}^\infty$ in such a way that $j\mapsto \Vert \kappa(j) \Vert_{\infty}$ is non-decreasing and consider functions on $[0,2\pi]^{d_U}$ of the form
\[
u = A \sum_{j=1}^\infty j^{-\eta} y_j e_j,
\qquad y_j\in [0,1].
\]


First, since $\|e_j\|_{\Lp{\infty}(\Omega_U)} = c_r^{-1}$,
\(
\|u\|_{\Lp{\infty}(\Omega_U)}
\le
Ac_r^{-1}\sum_{j=1}^\infty j^{-\eta}
\)
and the latter converges because $\eta>1$. Choosing $A(\eta)>0$ sufficiently small furthermore ensures that
\(
\|u\|_{\Lp{\infty}(\Omega_U)} \leq \beta_U
\)
uniformly. 

Second, we estimate the Lipschitz constant. For a function $f:[0,2\pi]^{d_U} \mapsto \bbR$, we define \[
\textrm{Lip}(f) = \inf \{L > 0 \,\mid \, \vert f(x) - f(y) \vert \leq L \Vert x - y \Vert_2 \text{ for all $x,y$}\}
\]
and note that that $\textrm{Lip}(f_1 + f_2)\leq \textrm{Lip}(f_1) + \textrm{Lip}(f_2)$ as well as $\textrm{Lip}(f) \leq C \Vert \nabla f \Vert_{\Lp{\infty}(\Omega_U)}$ by the mean-value theorem.
By noting that 
\(
\Lip(e_\kappa) \leq C \Vert \nabla e_\kappa \Vert_{\Lp{\infty}(\Omega_U)} \leq C \Vert \kappa \Vert_{\infty},
\)
we estimate as follows:
\begin{equation} \label{eq:hypercudeBoundedLipschitz:eq1}
\Lip(u)
\le
A \sum_{j=1}^\infty j^{-\eta}\Lip(e_j)
\leq C
A \sum_{j=1}^\infty j^{-\eta} \Vert \kappa(j) \Vert_{\infty}.  
\end{equation}
Let us now consider the inverse enumeration $\kappa \mapsto j(\kappa)$. For a fixed $\kappa \in \bbN^{d_U}$, we write $K = \Vert \kappa \Vert_\infty$ and, due to the fact that $j \mapsto \kappa(j)$ is non-decreasing, we have 
\begin{equation} \label{eq:hypercudeBoundedLipschitz:eq2}
  (K-1)^{d_U} \leq \vert \{\kappa' \in \bbN^{d_U} \, \mid \, \Vert \kappa' \Vert_\infty < K \} \vert \leq j(\kappa) \leq \vert \{\kappa' \in \bbN^{d_U} \, \mid \, \Vert \kappa' \Vert_\infty \leq K \} \vert = K^{d_U}
\end{equation}
as well as \begin{equation} \label{eq:hypercudeBoundedLipschitz:eq5}
    \vert \{\kappa' \in \bbN^{d_U} \, \mid \, \Vert \kappa' \Vert_\infty = K \} \vert \leq d_U K^{d_U -1},
\end{equation}
since fixing one coordinate equal to $K$ leaves at most $K^{d_U-1}$ choices for the remaining coordinates, and there are $d_U$ possible coordinates that may attain the maximum.
We continue our estimation: \begin{align}
\sum_{j=1}^\infty j^{-\eta} \Vert \kappa(j) \Vert_{\infty} &=  \sum_{K=1}^\infty \sum_{\Vert \kappa(j) \Vert_\infty = K} j(\kappa)^{-\eta} K \notag \\
&= \sum_{\Vert \kappa(j) \Vert_\infty = 1} j(\kappa)^{-\eta} + \sum_{K=2}^\infty \sum_{\Vert \kappa(j) \Vert_\infty = K} j(\kappa)^{-\eta} K \notag \\
&\leq1 + \sum_{K=2}^\infty \sum_{\Vert \kappa(j) \Vert_\infty = K} K^{1}(K-1)^{-\eta d_U} \label{eq:hypercudeBoundedLipschitz:eq3}\\
&\leq C \sum_{K=2}^\infty K^{d_U} (K-1)^{-\eta d_U}\label{eq:hypercudeBoundedLipschitz:eq4} \\
&\leq C \sum_{K=2}^\infty K^{d_U(1-\eta)}\label{eq:hypercudeBoundedLipschitz:eq6}
\end{align}
where we used \eqref{eq:hypercudeBoundedLipschitz:eq2} for \eqref{eq:hypercudeBoundedLipschitz:eq3}, \eqref{eq:hypercudeBoundedLipschitz:eq5} for \eqref{eq:hypercudeBoundedLipschitz:eq4} and the fact that $K/2 \leq K-1$ for $K \geq 2$ for \eqref{eq:hypercudeBoundedLipschitz:eq6}. The latter quantity converges since $\eta > 1 + \frac{1}{d_U}$ and, from \eqref{eq:hypercudeBoundedLipschitz:eq1}, we therefore deduce that $A(\eta)$ can be chosen so that $\textrm{Lip}(u) \leq L_U$ uniformly. This concludes the proof that $U$ contains an infinite-dimensional hypercube.
\end{proof}

\end{document}
